\documentclass[a4paper,11pt]{article}
\usepackage[T2A]{fontenc}
\usepackage[utf8]{inputenc}
\usepackage[russian]{babel}
\usepackage{amsmath,amssymb,mathtools}
\usepackage{geometry}
\geometry{left=23mm,right=23mm,top=22mm,bottom=23mm}
\usepackage{microtype}
\usepackage{icomma}
\usepackage{booktabs,tabularx,array}
\usepackage{enumitem}
\usepackage{xcolor}
\usepackage{tikz}
\usetikzlibrary{arrows.meta,positioning,shapes.geometric,calc}
\usepackage{hyperref}
\usepackage{fancyhdr}
\usepackage{setspace}
\usepackage{helvet}
\renewcommand{\familydefault}{\sfdefault}
\definecolor{accent}{HTML}{087D83}
\definecolor{darkaccent}{HTML}{07545B}
\definecolor{ink}{HTML}{303941}
\definecolor{mute}{HTML}{65727A}
\hypersetup{colorlinks=true,urlcolor=darkaccent,linkcolor=darkaccent}
\setlength{\parindent}{0pt}
\setlength{\parskip}{0.48em}
\setstretch{1.15}
\setlist[enumerate]{leftmargin=*,itemsep=.60em,topsep=.45em}
\pagestyle{fancy}
\fancyhf{}
\fancyhead[L]{\footnotesize\textcolor{mute}{Практико-ориентированные задачи}}
\fancyhead[R]{\footnotesize\textcolor{mute}{Анна Трапезникова}}
\fancyfoot[C]{\footnotesize\textcolor{mute}{\thepage}}
\renewcommand{\headrulewidth}{0pt}
\renewcommand{\footrulewidth}{0pt}
\setlength{\headheight}{14pt}
\newcommand{\topic}[1]{\section*{\textcolor{darkaccent}{#1}}}
\newcommand{\subtopic}[1]{\subsection*{\textcolor{accent}{#1}}}
\newcommand{\keyidea}[1]{\textcolor{darkaccent}{\textbf{Ключевая идея.}} #1}
\newcommand{\note}[1]{\textcolor{mute}{\textbf{Проверьте.}} #1}
\newcommand{\ex}[1]{\textbf{Пример.} #1}
\begin{document}
\thispagestyle{empty}
\begin{center}
\vspace*{25mm}
{\large\textcolor{darkaccent}{МАТЕМАТИКА}}\par
\vspace{24mm}
{\LARGE Чек-лист}\par
\vspace{16mm}
{\Large\color{accent} Практико-ориентированные задачи}\par
\vspace{3mm}
{\large\color{accent} Чтение условий, формулы и степени десяти}\par
\vspace{17mm}
{\normalsize Математика, 10 класс}\par
\vfill
{\large Анна Трапезникова}\par
\vspace{14mm}
{\normalsize 1 ноября 2026 г.}\par
\vspace{21mm}
{\small Лёвин Артём Александрович}\par
{\small эксклюзивно для Анны Трапезниковой}\par
\vspace{19mm}
\begin{tikzpicture}[x=1cm,y=1cm]
\draw[accent,line width=.8pt] (0,0) .. controls (1.5,.8) and (3.2,-.65) .. (5,0)
 .. controls (6.7,.6) and (8.5,-.7) .. (10,0);
\fill[accent] (5,0) circle (1.15pt);
\end{tikzpicture}
\end{center}
\clearpage

\topic{1. Главный метод: от текста к формуле}

\keyidea{Сначала установите, что означает каждая буква и какое значение ей соответствует. Затем подставляйте числа и выполняйте вычисления.}

\begin{enumerate}[label=\textcolor{accent}{\arabic*.},itemsep=.3em]
\item Прочитайте вопрос. Назовите искомую величину и её единицу измерения.
\item Найдите в условии формулу и составьте краткую таблицу обозначений.
\item Переведите данные в нужные единицы: миллиметры в метры, проценты в доли, большие числа в степени десяти.
\item Подставьте значения, сохраняя скобки и структуру дробей.
\item Выразите неизвестное, сократите общие множители и посчитайте.
\item Проверьте ответ: единицы измерения, знак, ограничения и смысл задачи.
\end{enumerate}

\subtopic{Что обязательно фиксировать}
\begin{center}
\begin{tabularx}{\linewidth}{@{}p{.28\linewidth}X@{}}
\toprule
\textbf{В условии} & \textbf{Что записать}\\
\midrule
Число с единицей & Например, $83\ \text{кг}$, $90^\circ\!\text{C}$ или $400\,000\ \text{Па}$.\\
Процент & $21\%=0{,}21$ как доля; если в формуле стоит $\cdot100\%$, учитывайте этот множитель отдельно.\\
Слово «не более» & Неравенство $\le$ (верхняя граница).\\
Слово «не менее» & Неравенство $\ge$ (нижняя граница).\\
Расстояние по горизонтали & Подставляют вместо $x$, если ось $Ox$ направлена горизонтально.\\
Высота или вертикальное смещение & Это значение $y$, если ось $Oy$ направлена вертикально.\\
\bottomrule
\end{tabularx}
\end{center}

\note{Встречающиеся рядом дополнительные пояснения об интервалах допустимых оценок не заменяют конкретных чисел, данных в условии.}

\subtopic{Первый образец: прочитать координатную модель}
Для цепи моста задана высота $y=0{,}005x^2-0{,}74x+25$, где $x$~--- расстояние от пилона вдоль моста. На расстоянии $30$ м подставляем \emph{именно $x=30$}:
\[
y(30)=0{,}005\cdot30^2-0{,}74\cdot30+25
=4{,}5-22{,}2+25=7{,}3\ \text{м}.
\]
Искомая высота равна $7{,}3$ м.

\clearpage
\topic{2. Степени десяти как инструмент вычисления}

\keyidea{Сдвигайте десятичную запятую и одновременно меняйте степень десяти так, чтобы значение числа сохранялось.}

\[
4{,}2\cdot10^3=42\cdot10^2, \qquad
8{,}3\cdot10^6=83\cdot10^5,
\qquad 90=9\cdot10^1.
\]
\[
10^a\cdot10^b=10^{a+b}, \qquad
\frac{10^a}{10^b}=10^{a-b},\qquad
(10^a)^b=10^{ab}.
\]
\[
0{,}005=5\cdot10^{-3},\qquad
0{,}74=74\cdot10^{-2}.
\]

\ex{Вычислите два слагаемых $0{,}005\cdot30^2$ и $0{,}74\cdot30$.}
\[
0{,}005\cdot30^2=(5\cdot10^{-3})(9\cdot10^2)
=45\cdot10^{-1}=4{,}5,
\]
\[
0{,}74\cdot30=(74\cdot10^{-2})(3\cdot10^1)
=222\cdot10^{-1}=22{,}2.
\]

\subtopic{Как не потерять степень}
\begin{itemize}[leftmargin=*,itemsep=.25em]
\item $4{,}2\cdot10^3$ и $42\cdot10^2$ обозначают одно число.
\item При переносе $10^{-20}$ из знаменателя в числитель показатель меняет знак:
\[
\frac{1}{16\cdot10^{-20}}=\frac{10^{20}}{16}.
\]
\item Сокращают только \emph{множители}, а не части суммы: $\frac{6x}{3}=2x$, но $\frac{6+x}{3}\ne2+x$.
\item В выражении $10^{-8}\cdot10^{20}$ получается $10^{12}$.
\end{itemize}

\subtopic{Пример с изменением единиц}
Если рельс удлинился на $3{,}2$ мм, перед подстановкой в формулу длины запишите:
\[
3{,}2\ \text{мм}=3{,}2\cdot10^{-3}\ \text{м}.
\]
В формуле $l(T)=l_0(1+\alpha T)$ при $l_0=8$ м и $\alpha=2\cdot10^{-5}\ ^\circ\!\text{C}^{-1}$ получаем
\[
T=\frac{l-l_0}{l_0\alpha}
=\frac{3{,}2\cdot10^{-3}}{8\cdot2\cdot10^{-5}}
=20^\circ\!\text{C}.
\]

\clearpage
\topic{3. Как выразить неизвестную величину}

\subtopic{Неизвестная в знаменателе}
Если $A=\dfrac{B}{Cx}$, где $A,C\ne0$, то
\[
ACx=B \quad\Longrightarrow\quad x=\frac{B}{AC}.
\]
Умножение на $x$ и деление на $AC$ выполняются для \emph{всего} равенства.

\subtopic{Неизвестная под квадратом}
Если $D^2=A$ и по смыслу $D>0$, то $D=\sqrt A$; отрицательную ветвь отбрасываем.

\ex{Давление под цилиндрической колонной задаётся формулой}
\[
P=\frac{4mg}{\pi D^2}.
\]
При $m=1200$ кг, $g=10\ \text{м/с}^2$, $\pi=3$ и требовании $P\le400\,000$ Па:
\[
D^2\ge\frac{4mg}{\pi P_{\max}}
=\frac{4\cdot1200\cdot10}{3\cdot400\,000}
=0{,}04.
\]
\[
\boxed{D_{\min}=\sqrt{0{,}04}=0{,}2\ \text{м}.}
\]
Чем больше диаметр опоры, тем меньше давление. Поэтому в ответе нужно указать \emph{минимальный} допустимый диаметр.

\subtopic{Неизвестная в квадратном уравнении}
В модели момента инерции катушки
\[
I=\frac{(m+2M)R^2}{2}+M(2Rh+h^2)
\]
после подстановки $m=6$, $M=2$, $R=6$, $I=270$ имеем
\[
270=\frac{(6+2\cdot2)\cdot36}{2}+2(12h+h^2)
=180+24h+2h^2.
\]
\[
2h^2+24h-90=0\quad\Longrightarrow\quad
h^2+12h-45=0.
\]
Корни: $h=3$ и $h=-15$. Расстояние $h$ положительно, поэтому \textbf{$h=3$}. Единица для $h$ та же, что для радиуса $R$.

\note{При возведении в квадрат, извлечении корня и решении квадратного уравнения обязательно проверяйте физический смысл корня.}

\clearpage
\topic{4. Проценты и формулы с КПД}

КПД можно записать двумя равноценными способами:
\[
\eta_{\text{доля}}=\frac{Q_{\text{полезн}}}{Q_{\text{затрач}}},
\qquad
\eta_{\%}=\frac{Q_{\text{полезн}}}{Q_{\text{затрач}}}\cdot100\%.
\]
Следовательно, $21\%=0{,}21$ как доля. Подставлять число $21$ вместо $0{,}21$ в первой формуле \emph{нельзя}.

\subtopic{Разобранная задача: сколько потребуется дров}
При нагревании $83$ кг воды с $10^\circ\text{C}$ до $100^\circ\text{C}$ заданы
\[
c=4{,}2\cdot10^3\ \frac{\text{Дж}}{\text{кг}\cdot{}^\circ\!\text{C}},
\qquad q=8{,}3\cdot10^6\ \frac{\text{Дж}}{\text{кг}},
\qquad\eta\le21\%.
\]
Количество теплоты для нагрева воды: $Q_{\text{полезн}}=cm_{\text{в}}(t_2-t_1)$. Для дров: $Q_{\text{затрач}}=qm_{\text{др}}$.

Минимальная масса дров достигается при максимально допустимом КПД, то есть при $\eta=0{,}21$:
\[
0{,}21=\frac{(4{,}2\cdot10^3)\cdot83\cdot90}
{(8{,}3\cdot10^6)m_{\text{др}}}.
\]
Представим $4{,}2\cdot10^3=42\cdot10^2$ и $8{,}3\cdot10^6=83\cdot10^5$, сократим число $83$:
\[
m_{\text{др}}
=\frac{42\cdot10^2\cdot83\cdot9\cdot10^1}{83\cdot10^5\cdot(21/100)}
=\frac{42\cdot9\cdot10^3}{21\cdot10^3}
=18\ \text{кг}.
\]
\textbf{Ответ: $18$ кг.}

\note{Если ошибочно поставить в формулу КПД число $21$ вместо $0{,}21$, результат уменьшится ровно в $100$ раз. Это типичная ловушка при расчётах с процентами.}

\subtopic{Минимум и максимум в тексте задачи}
Если КПД \emph{не более} $21\%$, то $\eta\le0{,}21$. Для получения заданного количества тепла при меньшем КПД нужно больше топлива. Значит, минимальную массу топлива ищут при $\eta=0{,}21$.

\clearpage
\topic{5. Сложные формулы: общая техника остаётся той же}

\subtopic{Четвёртая степень и закон Стефана--Больцмана}
По условию задачи дана формула $P=\sigma S T^4$. Чтобы найти температуру,
\[
T^4=\frac{P}{\sigma S},\qquad T=\sqrt[4]{\frac{P}{\sigma S}},\quad T>0.
\]
При $\sigma=5{,}7\cdot10^{-8}$, $S=\dfrac{10^{20}}{16}$ и $P=9{,}12\cdot10^{25}$:
\[
T^4=
\frac{9{,}12\cdot10^{25}\cdot16}{5{,}7\cdot10^{12}}
=25{,}6\cdot10^{13}
=2{,}56\cdot10^{14}.
\]
Так как $4^4=256$ и $(10^3)^4=10^{12}$, получаем
\[
T^4=256\cdot10^{12}=(4\cdot10^3)^4,
\qquad \boxed{T=4000\ \text{К}}.
\]

\subtopic{Рейтинг магазина: сначала подставить все данные}
Рейтинг вычисляют по формуле
\[
R=r_{\text{пок}}-
\frac{r_{\text{пок}}-r_{\text{экс}}}
{(K+1)\,\dfrac{0{,}02K}{r_{\text{пок}}+0{,}1}}.
\]
Здесь $K$~--- число отзывов, $r_{\text{пок}}$~--- средняя оценка покупателей, $r_{\text{экс}}$~--- оценка экспертов. При $K=20$, $r_{\text{пок}}=0{,}65$, $r_{\text{экс}}=0{,}37$:
\[
(K+1)\frac{0{,}02K}{r_{\text{пок}}+0{,}1}
=21\cdot\frac{0{,}4}{0{,}75}=11{,}2,
\]
\[
R=0{,}65-\frac{0{,}28}{11{,}2}
=0{,}65-0{,}025=\boxed{0{,}625}.
\]
Вложенная дробь~--- обычная дробь: вычисляйте её знаменатель по шагам.

\subtopic{Короткие модели из задач занятия}
\textbf{Прибыль.} Если $p$~--- цена, $v$~--- переменные затраты на единицу, $F$~--- постоянные расходы, $n$~--- выпуск, то
\[
\Pi=n(p-v)-F.
\]
Для требования «прибыль не менее $\Pi_0$» имеем $n\ge\dfrac{\Pi_0+F}{p-v}$ при $p>v$. Если $n$ измеряется штуками, округляйте вверх.

\textbf{Подъём уровня воды.} Если время падения уменьшилось с $t_1$ до $t_2$ и $h=5t^2$, то высота подъёма равна
\[
\Delta h=5t_1^2-5t_2^2
=5(t_1-t_2)(t_1+t_2).
\]

\clearpage
\thispagestyle{plain}
\begin{center}
{\Large\textcolor{darkaccent}{Алгоритм: решить задачу по готовой формуле}}
\end{center}
\vspace{10mm}
\begin{center}
\tikzset{
 start/.style={ellipse,draw=accent,thin,align=center,minimum width=5.5cm,minimum height=1.0cm,font=\small},
 action/.style={rounded corners=3pt,draw=accent!75,thin,align=center,text width=7.5cm,minimum height=1.1cm,font=\small,inner sep=5pt},
 decision/.style={diamond,draw=accent!80,thin,align=center,aspect=2.2,text width=4.3cm,inner sep=4pt,font=\small},
 arrow/.style={-{Stealth[length=2.5mm]},line width=.7pt,draw=darkaccent}
}
\begin{tikzpicture}[node distance=11mm and 18mm]
\node[start] (a) {Прочитать вопрос и найти искомое};
\node[action,below=of a] (b) {Соотнести обозначения формулы с данными условия};
\node[decision,below=of b] (c) {Единицы и проценты приведены?};
\node[action,right=21mm of c,text width=3.15cm] (d) {Перевести единицы и проценты в доли};
\node[action,below=15mm of c] (e) {Подставить данные и выразить неизвестное};
\node[action,below=of e] (f) {Упростить степени десяти и сократить множители};
\node[decision,below=of f] (g) {Ответ удовлетворяет ограничениям?};
\node[action,right=21mm of g,text width=3.15cm] (h) {Проверить подстановку, арифметику, знаки};
\node[start,below=16mm of g] (i) {Ответ с единицей измерения};
\draw[arrow] (a)--(b);\draw[arrow] (b)--(c);
\draw[arrow] (c)-- node[left,font=\scriptsize]{да} (e);
\draw[arrow] (c)-- node[above,font=\scriptsize]{нет} (d);
\draw[arrow] (d.north) |- ([xshift=6mm,yshift=6mm]c.north) -| (c.north);
\draw[arrow] (e)--(f);\draw[arrow] (f)--(g);
\draw[arrow] (g)--node[left,font=\scriptsize]{да}(i);
\draw[arrow] (g)--node[above,font=\scriptsize]{нет}(h);
\draw[arrow] (h.north) |- (e.east);
\end{tikzpicture}
\end{center}
\clearpage

\topic{Тренировочный блок --- Путь А}
\textcolor{mute}{10 задач от базовых вычислений до комплексных формул. В ответах указывайте единицы измерения.}

\begin{enumerate}[label=\textbf{\arabic*.},itemsep=1.0em]
\item Вычислите, используя степени десяти:
\[
\frac{(4{,}2\cdot10^3)(9\cdot10^1)}{7\cdot10^4}.
\]

\item Высота цепи моста описывается формулой $y=0{,}004x^2-0{,}5x+21$, где $x$~--- расстояние вдоль моста от пилона (в метрах). Найдите высоту цепи на расстоянии $25$ м от пилона.

\item До дождя камень падал до поверхности воды $1$ с. После дождя время падения уменьшилось на $0{,}3$ с. На сколько метров поднялась вода? Используйте $h=5t^2$.

\item Рельс при $0^\circ\text{C}$ имеет длину $12$ м. При нагревании он удлинился на $5{,}4$ мм. Найдите температуру, если $l(T)=l_0(1+\alpha T)$, где $\alpha=1{,}5\cdot10^{-5}\ ^\circ\!\text{C}^{-1}$.

\item Товар продают по $1000$ руб. за штуку. Переменные затраты~--- $650$ руб. на штуку, постоянные расходы~--- $420\,000$ руб. Найдите наименьший целый объём выпуска, обеспечивающий прибыль не менее $980\,000$ руб.

\item На цилиндрическую колонну опирается конструкция массой $900$ кг. Давление вычисляют по формуле $P=\dfrac{4mg}{\pi D^2}$. Найдите минимальный диаметр $D$, при котором $P\le300\,000$ Па. Примите $g=10\ \text{м/с}^2$, $\pi=3$.

\item Для рейтинга магазина используется формула из раздела 5. Найдите $R$, если $K=9$, $r_{\text{пок}}=0{,}8$, $r_{\text{экс}}=0{,}6$.

\item Для катушки $I=\dfrac{(m+2M)R^2}{2}+M(2Rh+h^2)$. Найдите $h\ge0$, если $m=4$, $M=1$, $R=5$, $I=131$ (в согласованных единицах).

\item КПД кормозапарника не превышает $30\%$. Какова \emph{наименьшая масса} дров (в килограммах), необходимая для нагревания $100$ кг воды от $20^\circ\text{C}$ до $80^\circ\text{C}$? Используйте
\[
\eta=\frac{cm_{\text{в}}(t_2-t_1)}{qm_{\text{др}}},\quad
c=4{,}2\cdot10^3,\quad q=7\cdot10^6,
\]
где $c$ и $q$ заданы в соответствующих единицах СИ, а $\eta$ записывается долей.

\item Мощность излучения звезды равна $P=5{,}7\cdot10^{24}$ Вт, площадь поверхности $S=\dfrac{10^{20}}{16}\ \text{м}^2$. Найдите температуру $T>0$ по закону $P=\sigma ST^4$, где $\sigma=5{,}7\cdot10^{-8}\ \text{Вт}/(\text{м}^2\cdot\text{К}^4)$.
\end{enumerate}

\clearpage
\topic{Ответы к тренировочному блоку}
\vspace{5mm}
\begin{center}
\renewcommand{\arraystretch}{1.55}
\begin{tabularx}{.84\linewidth}{@{}cX@{}}
\toprule
\textbf{№} & \textbf{Ответ}\\
\midrule
1 & $5{,}4$\\
2 & $11$ м\\
3 & $2{,}55$ м\\
4 & $30^\circ\text{C}$\\
5 & $4000$ штук\\
6 & $0{,}2$ м\\
7 & $0{,}7$\\
8 & $h=4$ (в единицах радиуса $R$)\\
9 & $12$ кг\\
10 & $2000$ К\\
\bottomrule
\end{tabularx}
\end{center}
\vspace{12mm}
\topic{Итоговая памятка}
\begin{center}
\begin{tabularx}{\linewidth}{@{}p{.34\linewidth}X@{}}
\toprule
\textbf{Сигнал в задаче} & \textbf{Полезное действие}\\
\midrule
Громоздкая формула & Расшифровать каждую букву до подстановки.\\
Много нулей и запятых & Использовать степени десяти и сокращения.\\
КПД в процентах & Перевести процент в долю либо сохранить множитель $100\%$.\\
«Не менее», «не более» & Составить неравенство и проверить граничное значение.\\
Квадрат или четвёртая степень & Найти допустимый корень и проверить его смысл.\\
График и расстояние & Сначала определить, какой координате соответствует данное расстояние.\\
\bottomrule
\end{tabularx}
\end{center}
\vfill
\begin{center}
{\small\color{mute}Лёвин Артём Александрович эксклюзивно для Анны Трапезниковой}
\end{center}
\end{document}
